\documentclass{article}
\usepackage[portuguese]{babel}
\usepackage{listings}
\usepackage{xcolor}
\usepackage[margin=1.5cm]{geometry}

\definecolor{codegreen}{rgb}{0,0.6,0}
\definecolor{codegray}{rgb}{0.5,0.5,0.5}
\definecolor{codepurple}{rgb}{0.8,0,0.2}
\definecolor{backcolour}{rgb}{.8,.8,1}

\lstdefinestyle{mystyle}{
    backgroundcolor=\color{backcolour},   
    commentstyle=\color{codegreen},
    keywordstyle=\color{blue},
    numberstyle=\tiny\color{codegray},
    stringstyle=\color{codepurple},
    basicstyle=\ttfamily\footnotesize,
    breakatwhitespace=false,         
    breaklines=true,                 
    captionpos=b,                    
    keepspaces=true,                 
    numbers=left,                    
    numbersep=5pt,                  
    showspaces=false,                
    showstringspaces=false,
    showtabs=false,                  
    tabsize=2
}

\lstset{style=mystyle}
\begin{document}

\section{AA9}
Escreva um programa, em linguagem Python, para resolver um sistema de equações lineares triangular superior a partir das seguintes etapas:
\begin{enumerate}
    \item ler a dimensão n da matriz dos coeficientes no console;
    \item registrar as entadas da matriz dos coeficientes no console;
    \item registrar as entradas do vetor dos resultados no console;
    \item imprimir o sistema original na forma de matriz ampliada;
    \item resolver e imprimir a solução do sistema de equações;
\end{enumerate}
\subsection{}

Os item 1,2 e 3 desse exercicio ja foi resolvido no AA8. Sendo assim partiremos diretamente para os item 4 e 5.Para o item 4, definiremos uma funcao no inicio do nosso codigo assim como fizmos antes.

\begin{lstlisting}[language=python]
def printAmpliada(matriz, vetor):
    for i in range(len(matriz)):
        print(f"{matriz[i]}{vetor[i]}")
\end{lstlisting}

\subsection{Solucao para sistema triangular superior}
Para resolver esse tipo de sistema, obtemos primeiro a solucao do ultimo elemento e vamos substituindo os resultados de traz pra frente.

\begin{lstlisting}[language=Python]
incognitas = [0]*ordem
    incognitas[ordem-1] = vetorResultado[ordem-1]/matriz[ordem-1][ordem-1]
    for n in range(1,ordem):
        somatorio = 0
        k = ordem - n-1
        for i in range(n):
            j = ordem - i - 1
            somatorio += matriz[k][j] * incognitas[j]
        incognitas[k] = (vetorResultado[k]-somatorio)/matriz[k][k]

    print(incognitas)
\end{lstlisting}

Por estarmos trabalhando de traz pra frente, o algoritmo fica mais dificil de entender. Por isso recomendo fazer primeiro o AA10, e depois voltar

\subsection{Codigo completo:}
\begin{lstlisting}[language=Python]
deveParar = ""

def printmatriz(matriz):
    for linha in matriz:
        print(linha)

def printAmpliada(matriz, vetor):
    for i in range(len(matriz)):
        print(f"{matriz[i]}{vetor[i]}")



while (deveParar == ""):
    
    ordem = int(input("Digite a dimensao da matriz A\n"))
    matriz =[]
    for u in range(ordem):
        linha = []
        for v in range(ordem):
            linha.append(int(input(f"Digite o valor na posisao {u+1},{v+1} \n")))
        matriz.append(linha)
    print("Matriz inicial A:")
    printmatriz(matriz)
    vetorResultado = []
    for i in range(ordem):
        vetorResultado.append(int(input(f"Digite o valor de X{i} do vetor resultado")))
    print(vetorResultado)


    printmatriz(matriz)
    printAmpliada(matriz,vetorResultado)

    incognitas = [0]*ordem
    incognitas[ordem-1] = vetorResultado[ordem-1]/matriz[ordem-1][ordem-1]
    for n in range(1,ordem):
        somatorio = 0
        k = ordem - n-1
        for i in range(n):
            j = ordem - i - 1
            somatorio += matriz[k][j] * incognitas[j]
        incognitas[k] = (vetorResultado[k]-somatorio)/matriz[k][k]

    print(incognitas)
    deveParar = input("Deseja continuar?(enter para continuar, qualquer tecla para parar)\n")

print("Ate mais")
\end{lstlisting}

\end{document}
